\clearpage
\begingroup
\fancyhead[L]{\small THEORIE DER BEWEGUNG DES MONDES.}
\renewcommand{\headrulewidth}{0pt}
\par\noindent Hieraus folgt
\[
\begin{aligned}
0&=\frac{\mathrm{dd}(x-x')}{\mathrm{d}t^2}
  +\frac{(m+m')(x-x')}{R''^3}\\
 &\qquad+m''\left\{\frac{x-x'}{R'^3}
  -(x''-x')\left(\frac{1}{R'^3}-\frac{1}{R^3}\right)\right\}\\[5pt]
0&=\frac{\mathrm{dd}(y-y')}{\mathrm{d}t^2}
  +\frac{(m+m')(y-y')}{R''^3}\\
 &\qquad+m''\left\{\frac{y-y'}{R'^3}
  -(y''-y')\left(\frac{1}{R'^3}-\frac{1}{R^3}\right)\right\}\\[5pt]
0&=\frac{\mathrm{dd}(z-z')}{\mathrm{d}t^2}
  +\frac{(m+m')(z-z')}{R''^3}\\
 &\qquad+m''\left\{\frac{z-z'}{R'^3}
  -(z''-z')\left(\frac{1}{R'^3}-\frac{1}{R^3}\right)\right\}
  \mathpunct{\raisebox{.4ex}{.}}
\end{aligned}
\]

\par\noindent\hspace*{\parindent}Endlich bezeichnen wir Länge des Mondes und der Sonne, vom Mittelpunkt
der Erde gesehen und vom Frühlingsnachtgleichepunkt zur Zeit \(T\)
gezählt, durch \(v\), \(V\), Breite des Mondes durch \(\beta\), und setzen
\(\operatorname{tang}\beta=\theta\), \(R''\cos\beta=r\), \(\frac{1}{r}=p\). Aus
diesen Voraussetzungen folgt leicht
\[
\begin{array}{lll}
x-x'=r\cos v=\dfrac{\cos v}{p},&
y-y'=\dfrac{\sin v}{p},&z-z'=\dfrac{\theta}{p}\\[5pt]
x''-x'=R\cos V,&y''-y'=R\sin V,&
\end{array}
\]
\noindent\(z''-z'=0\) (weil wir die Breite der Sonne \(=0\) voraussetzen).

\par\noindent\hspace*{\parindent}Durch Substitution dieser Werthe in obigen Gleichungen erhalten wir
folgende neue:
\[
\begin{aligned}
0&=-\frac{\cos v\, .\,\mathrm{dd}p}{pp\,\mathrm{d}t^2}
  +\frac{2\cos v\, .\,\mathrm{d}p^2}{p^3\mathrm{d}t^2}
  +\frac{2\sin v\, .\,\mathrm{d}p\mathrm{d}v}{pp\,\mathrm{d}t^2}\\
 &\qquad-\frac{\sin v\, .\,\mathrm{dd}v}{p\,\mathrm{d}t^2}
  -\frac{\cos v\, .\,\mathrm{d}v^2}{p\,\mathrm{d}t^2}
  +\frac{(m+m')pp\cos v}{(1+\theta\theta)^{\frac{3}{2}}}\\
 &\qquad+m''\left\{\frac{\cos v}{pR'^3}
  -R\cos V\left(\frac{1}{R'^3}-\frac{1}{R^3}\right)\right\}\\[6pt]
0&=-\frac{\sin v\, .\,\mathrm{dd}p}{pp\,\mathrm{d}t^2}
  +\frac{2\sin v\, .\,\mathrm{d}p^2}{p^3\mathrm{d}t^2}
  -\frac{2\cos v\, .\,\mathrm{d}p\mathrm{d}v}{pp\,\mathrm{d}t^2}\\
 &\qquad+\frac{\cos v\, .\,\mathrm{dd}v}{p\,\mathrm{d}t^2}
  -\frac{\sin v\, .\,\mathrm{d}v^2}{p\,\mathrm{d}t^2}
  +\frac{(m+m')pp\sin v}{(1+\theta\theta)^{\frac{3}{2}}}\\
 &\qquad+m''\left\{\frac{\sin v}{pR'^3}
  -R\sin V\left(\frac{1}{R'^3}-\frac{1}{R^3}\right)\right\}\\[6pt]
0&=-\frac{\theta\mathrm{dd}p}{pp\,\mathrm{d}t^2}
  +\frac{2\theta\mathrm{d}p^2}{p^3\mathrm{d}t^2}
  -\frac{2\mathrm{d}p\mathrm{d}\theta}{pp\,\mathrm{d}t^2}
  +\frac{\mathrm{dd}\theta}{p\,\mathrm{d}t^2}\\
 &\qquad+\frac{(m+m')pp\theta}{(1+\theta\theta)^{\frac{3}{2}}}
  +\frac{m''\theta}{pR'^3}\mathpunct{\raisebox{.4ex}{.}}
\end{aligned}
\]

\par\noindent\hspace*{\parindent}Die Combination dieser Gleichungen gibt uns folgende:
\[
\begin{aligned}
0&=-\frac{\mathrm{dd}p}{\mathrm{d}t^2}
  +\frac{2\mathrm{d}p^2}{p\,\mathrm{d}t^2}
  -\frac{p\,\mathrm{d}v^2}{\mathrm{d}t^2}
  +\frac{(m+m')p^4}{(1+\theta\theta)^{\frac{3}{2}}}\\
 &\qquad+m''\left\{\frac{p}{R'^3}
  -Rpp\cos(v-V)\left(\frac{1}{R'^3}-\frac{1}{R^3}\right)\right\}
\end{aligned}
\]
\SourceLeaf{622}
\[
\begin{aligned}
0&=-\frac{2\mathrm{d}p\mathrm{d}v}{p\,\mathrm{d}t^2}
  +\frac{\mathrm{dd}v}{\mathrm{d}t^2}
  +m''Rp\sin(v-V)\left(\frac{1}{R'^3}-\frac{1}{R^3}\right)\\[5pt]
0&=\frac{\mathrm{dd}\theta}{\mathrm{d}t^2}
  +\frac{\theta\mathrm{d}v^2}{\mathrm{d}t^2}
  -\frac{2\mathrm{d}p\mathrm{d}\theta}{p\,\mathrm{d}t^2}\\
 &\qquad+m''\theta Rp\cos(v-V)\left(\frac{1}{R'^3}-\frac{1}{R^3}\right)
  \mathpunct{\raisebox{.4ex}{.}}
\end{aligned}
\]

\par\noindent Setzen wir hier wiederum
\(\dfrac{\mathrm{dd}p}{\mathrm{d}t^2}
 =\dfrac{\mathrm{d}v^2}{\mathrm{d}t^2}
 \left(\dfrac{\mathrm{dd}p}{\mathrm{d}v^2}
 -\dfrac{\mathrm{dd}t}{\mathrm{d}v^2}\dfrac{\mathrm{d}p}{\mathrm{d}t}\right)\),
\(\dfrac{\mathrm{dd}v}{\mathrm{d}t^2}
 =-\dfrac{\mathrm{d}v^3}{\mathrm{d}t^3}
 \dfrac{\mathrm{dd}t}{\mathrm{d}v^2}\),
\(\dfrac{\mathrm{dd}\theta}{\mathrm{d}t^2}
 =\dfrac{\mathrm{d}v^2}{\mathrm{d}t^2}
 \left(\dfrac{\mathrm{dd}\theta}{\mathrm{d}v^2}
 -\dfrac{\mathrm{dd}t}{\mathrm{d}v^2}\dfrac{\mathrm{d}\theta}{\mathrm{d}t}\right)\),
nach Gründen der Integralrechnung, und Kürze halber die mittlere Bewegung
des Mondes in der Länge \(=u\), und also
\(\dfrac{\mathrm{d}u}{\mathrm{d}t}=\mathit{Const.}\), [sowie]
\[
\begin{aligned}
m''\frac{\mathrm{d}t^2}{\mathrm{d}u^2}Rp
 \left(\frac{1}{R'^3}-\frac{1}{R^3}\right)\sin 2(v-V)&=\pi\\[5pt]
m''\frac{\mathrm{d}t^2}{\mathrm{d}u^2}Rp
 \left(\frac{1}{R'^3}-\frac{1}{R^3}\right)\cos 2(v-V)&=\omega\\[5pt]
\omega-m''\frac{\mathrm{d}t^2}{\mathrm{d}u^2}\frac{1}{R'^3}&=\psi,
\end{aligned}
\]
\noindent so werden unsere Gleichungen:
\[
\begin{aligned}
0&=\frac{\mathrm{dd}p}{\mathrm{d}v^2}
  -\frac{\mathrm{dd}u}{\mathrm{d}v^2}\frac{\mathrm{d}p}{\mathrm{d}u}
  -\frac{2\mathrm{d}p^2}{p\,\mathrm{d}v^2}+p\\
 &\qquad-\frac{p^4(m+m')}{(1+\theta\theta)^{\frac{3}{2}}}
  \frac{\mathrm{d}t^2}{\mathrm{d}u^2}\frac{\mathrm{d}u^2}{\mathrm{d}v^2}
  +p\psi\frac{\mathrm{d}u^2}{\mathrm{d}v^2}\\[5pt]
0&=\frac{2\mathrm{d}p}{p\,\mathrm{d}v}
  +\frac{\mathrm{dd}u}{\mathrm{d}v^2}\frac{\mathrm{d}v}{\mathrm{d}u}
  -\frac{\mathrm{d}u^2}{\mathrm{d}v^2}\pi\\[5pt]
0&=\frac{\mathrm{dd}\theta}{\mathrm{d}v^2}
  -\frac{\mathrm{dd}u}{\mathrm{d}v^2}\frac{\mathrm{d}\theta}{\mathrm{d}u}
  -\frac{2\mathrm{d}p\mathrm{d}\theta}{p\,\mathrm{d}v^2}
  +\theta+\theta\frac{\mathrm{d}u^2}{\mathrm{d}v^2}\omega.
\end{aligned}
\]

\par\noindent Schaffen wir endlich aus der ersten und dritten Gleichung, mit Hülfe der
zweiten, \(\dfrac{\mathrm{dd}u}{\mathrm{d}v^2}\) weg, so ergeben sich folgende drei
\textls[100]{Fundamentalgleichungen}:
\[
\begin{aligned}
\mathrm{I.}\qquad0&=\frac{\mathrm{dd}p}{\mathrm{d}v^2}+p
 -p^4\frac{\mathrm{d}u^2}{\mathrm{d}v^2}
 \left\{\frac{(m+m')}{(1+\theta\theta)^{\frac{3}{2}}}
 \frac{\mathrm{d}t^2}{\mathrm{d}u^2}\right.\\
 &\hspace{9em}\left.-\frac{\psi}{p^3}
 +\frac{\mathrm{d}p}{p^4\mathrm{d}v}\pi\right\}\\[6pt]
\mathrm{II.}\qquad0&=\frac{\mathrm{dd}u}{\mathrm{d}v^2}\frac{\mathrm{d}v}{\mathrm{d}u}
 +\frac{2\mathrm{d}p}{p\,\mathrm{d}v}
 -\frac{\mathrm{d}u^2}{\mathrm{d}v^2}\pi\\[6pt]
\mathrm{III.}\qquad0&=\frac{\mathrm{dd}\theta}{\mathrm{d}v^2}+\theta
 +\left(\theta\omega-\frac{\mathrm{d}\theta}{\mathrm{d}v}\pi\right)
 \frac{\mathrm{d}u^2}{\mathrm{d}v^2}\mathpunct{\raisebox{.4ex}{.}}
\end{aligned}
\]
\begin{center}\rule{.9in}{.35pt}\end{center}
\clearpage
\endgroup
